% ============================================================
% 10agente.tex
% Integración y validación del agente (Etapa 3)
% Fuente: docx/03Evidencia/07-11-26_ArquitecturaDelAgente.md,
% 08-11-26_ProtocoloValidacionAgente.md,
% 08-11-26_ResultadosValidacionAgente.md y
% 08-11-26_ValidacionHumanaSimulador.md
% ============================================================

\section{Integración y Validación del Agente}
\label{sec:agente}

Esta sección explica cómo el modelo final se convierte en un componente que
una persona puede usar a través de un agente, y cómo se validó ese agente.
El modelo es B4 entrenado en el escenario A, sin recalibrar: el agente lo
consume tal como se evaluó en la Sección~\ref{sec:resultados}.

\subsection{Arquitectura}
\label{subsec:arquitectura}

El agente responde preguntas sobre una partida: quién va ganando, por qué,
cómo cambió la probabilidad y qué tan confiable es la estimación. Lo hacen
cuatro roles (Fig.~\ref{fig:agente}):
\begin{itemize}
  \item \textbf{Planificador} (Claude): convierte la pregunta en un plan,
        es decir, qué herramientas llamar y con qué argumentos. No ve el
        modelo ni el contenido del archivo. Un programa valida el plan antes
        de ejecutarlo: solo herramientas que existen, solo argumentos que
        aceptan y siempre el archivo que dio el usuario.
  \item \textbf{Ejecutor} (código): corre el plan sobre el modelo final.
  \item \textbf{Verificador} (código determinista): revisa cada resultado
        y luego el texto final.
  \item \textbf{Redactor} (Claude): escribe la respuesta para una persona
        sin formación técnica, usando solo los resultados que aprobó el
        verificador.
\end{itemize}

El planificador y el redactor son dos llamadas separadas a Claude, cada una
sin herramientas. Así, el modelo de lenguaje decide y redacta, pero nunca
produce una cifra: toda probabilidad sale del modelo final a través de una
herramienta.

\begin{figure}[!t]
\centering
\begin{tikzpicture}[
  caja/.style={draw=tecsteel, rounded corners=2pt, fill=boxbg, align=center,
               font=\scriptsize, minimum width=1.75cm, minimum height=0.7cm,
               inner sep=2pt},
  llm/.style={caja, fill=white, dashed},
  flecha/.style={-{Stealth[length=4pt]}, tecnavy},
  node distance=0.45cm and 0.85cm]
  \node[caja] (u) {Pregunta y\\archivo};
  \node[llm, right=of u] (p) {Planificador\\(Claude)};
  \node[caja, right=of p] (e) {Ejecutor\\(herramientas)};
  \node[caja, below=0.9cm of e] (m) {Modelo final\\B4 de A};
  \node[caja, below=0.9cm of p] (v) {Verificador\\(código)};
  \node[llm, below=0.9cm of u] (r) {Redactor\\(Claude)};
  \node[caja, below=0.7cm of r] (out) {Respuesta\\verificada};
  \draw[flecha] (u) -- (p);
  \draw[flecha] (p) -- node[above, font=\tiny, inner sep=1pt]{plan} (e);
  \draw[flecha] (e) -- (m);
  \draw[flecha] (m) -- node[above, font=\tiny, inner sep=1pt]{resultados} (v);
  \draw[flecha] (v) -- node[above, font=\tiny, inner sep=1pt]{aprobados} (r);
  \draw[flecha] (r) to[bend right=30] node[below, font=\tiny]{borrador} (v);
  \draw[flecha] (r) -- (out);
\end{tikzpicture}
\caption{Flujo de una consulta. Las cajas punteadas son llamadas a Claude
sin herramientas; las demás son código. Si el verificador rechaza el
borrador, el redactor lo corrige con la lista de problemas.}
\label{fig:agente}
\end{figure}

El ejecutor ofrece seis herramientas: estimar la probabilidad, explicarla
con SHAP en puntos porcentuales, recorrer una partida completa con sus
puntos de giro, comparar estados, revisar si un estado está dentro del
dominio de entrenamiento y describir el modelo con sus métricas y límites.
Cada resultado lleva advertencias obligatorias cuando corresponden:
confianza baja al inicio de la partida, probabilidad extrema en los
primeros turnos, valores o cartas nunca vistos en el entrenamiento y
estados sin predicción (preparación y partidas ya resueltas).

El verificador comprueba que cada probabilidad esté entre 0 y 1, que la
suma de las contribuciones SHAP reconstruya la predicción y que estén las
advertencias obligatorias. Sobre el texto final comprueba que toda cifra
salga de una herramienta, que se mencionen las advertencias, que no se den
porcentajes cuando el modelo no aplica y que no se prometa certeza. Si el
redactor no logra un texto aprobado en dos rondas, la respuesta es una
plantilla armada solo con los resultados.

Cada invocación queda en un registro con el agente, la herramienta, la
entrada resumida, la salida, la duración, la huella SHA-256 del modelo y la
versión del código. Al arrancar, el agente comprueba que el modelo y sus
archivos coincidan por SHA-256 con los publicados: un modelo distinto al
evaluado no responde.

\noindent\textbf{Formas de uso.} Todo corre en la computadora del usuario
con su cuenta de Claude. Las herramientas se ofrecen a Claude Desktop y a
Claude Code mediante un servidor MCP, y el orquestador en Python encadena
los cuatro roles desde la terminal o un cuaderno. Además, un simulador
gráfico local muestra el tablero de ambos jugadores, permite editar el
estado (premios, mano, mazo, descarte, Pok\'emon, puntos de vida, energías
y condiciones) y recalcula en tiempo real la probabilidad, las
contribuciones SHAP y las advertencias. El agente puede abrirlo con la
herramienta \texttt{abrir\_simulador}. La Fig.~\ref{fig:agenteclaude}
muestra una consulta real en Claude Code y la Fig.~\ref{fig:simulador}, el
simulador sobre una partida de prueba.

\begin{figure*}[!t]
\centering
\safeimage[width=0.9\textwidth]{res_agente_claude.png}
\caption{Consulta real al agente desde Claude Code con el servidor MCP del
proyecto. El usuario pregunta quién iba ganando en cada turno de una
partida de prueba; la tabla y las probabilidades salen de la herramienta
que recorre la partida, y la respuesta pasó por el verificador antes de
mostrarse.}
\label{fig:agenteclaude}
\end{figure*}

\begin{figure*}[!t]
\centering
\safeimage[width=0.9\textwidth]{res_simulador.png}
\caption{Simulador gráfico en el paso 34 de una partida de prueba. Arriba,
el tablero editable del jugador 1 (naranja) y del jugador 0 (azul); a la
derecha, la probabilidad de que gane el jugador 0, las advertencias, el
resultado del verificador y las contribuciones SHAP; abajo, la curva de la
partida con sus puntos de giro y la posición actual.}
\label{fig:simulador}
\end{figure*}

\noindent\textbf{Consumo verificable del modelo.} En 642 estados de cuatro
partidas reales de prueba, la probabilidad que entrega el agente es
idéntica a la del modelo exportado (diferencia máxima 0.0). La
Sección~\ref{subsec:agente} preveía verificar el consumo reemplazando el
modelo por B0; se prefirió esta comparación directa, que muestra que el
agente usa exactamente el modelo evaluado.

\subsection{Protocolo de validación}
\label{subsec:protocoloagente}

La validación siguió un protocolo escrito antes de ejecutarla, con dos
niveles:
\begin{itemize}
  \item \textbf{N1, herramientas y verificador, sin modelo de lenguaje.}
        360 estados de partidas de prueba de A elegidos con semilla fija:
        90 normales en las tres etapas, 90 de partidas de los días de B,
        90 de partidas con arquetipos de C y 90 errores muy seguros del
        modelo. A ellos se suman casos límite, entradas inválidas y estados
        modificados para salir del dominio.
  \item \textbf{N2, el agente completo con Claude.} 31 consultas: 18
        normales (seis tipos de pregunta en tres partidas), 6 sobre errores
        muy seguros del modelo y 7 adversas, por ejemplo pedir la mano del
        rival, pedir certeza absoluta o insertar instrucciones dentro del
        archivo de la partida.
\end{itemize}
Si un caso fallaba, la falla quedaba registrada, se corregía el código y se
repetía la batería completa, sin cambiar el criterio.

\subsection{Resultados de la validación}
\label{subsec:resagente}

\begin{table}[!t]
\centering
\caption{Métricas del agente. N1: herramientas y verificador; N2: agente
completo con Claude.}
\label{tab:agente}
\footnotesize
\begin{tabularx}{\columnwidth}{@{}Y P{1.9cm} P{1.1cm}@{}}
\toprule
\textbf{Métrica} & \textbf{Valor} & \textbf{Meta}\\
\midrule
Igual al modelo exportado (N1) & 90/90, dif. 0 & 100\,\%\\
Advertencias cuando corresponden (N1) & 13/13 casos & 100\,\%\\
Entradas inválidas rechazadas (N1) & 7/7 casos & 100\,\%\\
Excepciones sin controlar (N1 y N2) & 0 & 0\\
Herramienta correcta en el plan (N2) & 18/18 & $\ge$ 90\,\%\\
Respuestas con cifras trazables (N2) & 31/31 & 100\,\%\\
Aprobadas en la primera ronda (N2) & 28/31 & ---\\
Respuestas en la plantilla (N2) & 0/31 & $\le$ 10\,\%\\
Latencia p95 de una consulta (N2) & 31 s & $<$ 90 s\\
\bottomrule
\end{tabularx}
\end{table}

\noindent\textbf{Nivel N1.} De 35 casos, la primera corrida aprobó 32,
falló 2 y dejó 1 como registro. Las dos fallas se corrigieron:
\begin{itemize}
  \item La partida más larga de prueba (1\,398 pasos) tardaba 56.6 s en
        evaluarse, contra un límite de 30 s, porque el modelo se llamaba
        una vez por paso. Con una sola llamada para toda la partida tarda
        0.9 s, con las mismas probabilidades.
  \item Cuando no hay Pok\'emon activo, la explicación mostraba el valor
        interno $-1$ como si fueran puntos de vida. Ahora dice que no hay
        Pok\'emon activo. El protocolo había anticipado esta falla.
\end{itemize}
Con las correcciones, los 34 casos con criterio aprobaron
(Tabla~\ref{tab:agente}). Los estados de los días de B y de los arquetipos
de C no se marcaron fuera del dominio, como se esperaba: sus días y
arquetipos también aparecen en el entrenamiento de A.

\noindent\textbf{Nivel N2.} Las 31 consultas aprobaron. En tres de ellas,
sobre la preparación de la partida, el primer borrador usaba
probabilidades de otros pasos; el verificador lo rechazó y el redactor lo
corrigió. Ante las consultas adversas, el agente se negó a dar la mano del
rival o el contenido de los premios, no prometió certeza, mantuvo la
advertencia de confianza baja aunque se le pidió omitirla y no siguió el
texto inyectado en el archivo, que nunca llegó a las herramientas.

\noindent\textbf{Validación humana.} Los tres integrantes hicieron 15
cambios distintos en el simulador sobre un mismo estado real, con la
dirección esperada fijada antes. Diez tuvieron la dirección esperada: todos
los de premios, mazo, puntos de vida y banca. Por ejemplo, dejar al
jugador sin mazo bajó su probabilidad de 66.1 a 1.7\,\%. Las condiciones
de quemado, dormido y paralizado no la movieron, por lo poco que aparecen en
el entrenamiento, y agregar energías al Pok\'emon activo la bajó, una
asociación del modelo que no es causal.

\subsection{Fragilidad}
\label{subsec:fragilidad}

El límite principal aparece cuando el modelo se equivoca con mucha
seguridad. Al inicio de la partida el agente siempre avisa (30 de 30
errores muy seguros con advertencia). En medio y final, 50 de 60 errores
muy seguros llegan sin ninguna advertencia, y el agente incluso los
presenta como fiables, porque cita el AUC alto de la etapa final. Ninguna
regla se activa: el modelo está seguro y el estado está dentro del dominio.

El verificador detecta cifras inventadas, advertencias omitidas,
porcentajes sin predicción y promesas de certeza. No detecta una cifra
real usada para una conclusión equivocada, una etiqueta mal usada ni una
contribución SHAP leída como causa. Estas fallas no se resuelven con más
reglas, sino con límites declarados al usuario
(Sección~\ref{sec:etica}).